\chapter{Common concepts}\label{ch:concepts}

This chapter describes the design common to all classes of the library. The class chapters refer to
it and describe only what is specific to the class.

\section{Reference counting and copy on demand}\label{sec:concepts:refcount}

\subsection{Shared data}

An object of \code{NRVec}, \code{NRMat}, \code{NRSMat} and of all other classes holding large data
(\code{SparseMat}, \code{SparseSMat}, \code{Tensor}, \code{fourindex}, \code{Polynomial}, the permutation
classes, \code{bitvector}, \dots) consists of a small header (dimensions, a pointer to the data and a pointer
to a reference counter) and the data block. The copy constructor and the assignment operator between objects
of the same type copy only the header and increment the reference counter, so that both objects share the
same data:
\begin{lstlisting}
NRVec<double> a(1.,1000);   //1000 elements equal to 1
NRVec<double> b(a);         //no copy of the data, a.getcount()==2
NRVec<double> c;
c = a;                      //no copy, count is 3 now
\end{lstlisting}
The data are deallocated when the last object referring to them is destroyed or assigned a different
value. The current value of the counter is returned by \code{getcount()}.

Consequently, objects can be passed to and returned from functions by value at low cost, and expressions
like \code{x = A*x} are safe, since the result is computed into a new temporary object which is then
assigned to \code{x}.

\subsection{Copy on demand}

Before shared data are modified, the object must obtain its own copy of them. This is done by the member
function
\begin{lareference}
\lafunc{void copyonwrite(bool detachonly=false, bool deep=true)}{If the reference count is larger than one,
allocate a private copy of the data (copy them unless \code{detachonly}, which only allocates
uninitialized memory) and decrement the count of the shared block. For nested element types (see
section~\ref{sec:concepts:traits}) it is applied recursively to the elements when \code{deep} is true, also when the
outer object is not shared. For an object with count one it does nothing (apart from the recursion).}
\end{lareference}

The rules are:
\begin{enumerate}
\item \emph{Every member function and operator of the library which modifies the data of an object as a whole calls
	\code{copyonwrite()} itself} (\code{axpy}, \code{gemv}, \code{gemm}, \code{+=}, \code{*=}, assignment of a
	scalar, \code{randomize}, \code{storesubmatrix}, \code{resize}, \dots). Explicit calls before these
	functions are not needed.
\item \emph{Before writing to individual elements through \code{operator[]}, \code{operator()}, a raw pointer}
	(\code{operator T*()}, \code{&x[0]}, \dots) \emph{or a member function accessing individual elements} (such as
	\code{swap(i,j)} of \code{NRVec}, used in sorting loops) the user has to call \code{copyonwrite()} explicitly.
	In the first cases the library cannot distinguish a write access from a read access; element-level members do not call it
	by design, since a \code{copyonwrite()} inside inner loops over individual elements would be an unacceptable overhead:
\begin{lstlisting}
NRVec<double> b(a);   //b shares data with a
b.copyonwrite();      //b gets its own copy
b[0] = 5.;            //a is not changed
\end{lstlisting}
	Without the \code{copyonwrite()}, the assignment would silently change \code{a} as well.
	The rule applies to every object which might share its data: objects obtained by copying or assignment,
	objects received as function arguments by reference or pointer (they may share data with other copies in
	the caller), objects stored in other containers. A freshly constructed object (by a constructor taking
	dimensions) has count one, but calling \code{copyonwrite()} for it costs nothing.
\item \code{resize()} detaches an object from shared data: a shared object always gets newly allocated
	memory (with the old contents only if their preservation is requested), even when the size does not change.
	An object which has just been resized can be written to. \code{resize(0)} deallocates.
\end{enumerate}

\subsection{The DEBUG check of the reference count}

In programs compiled with \code{DEBUG}, the non-const versions of \code{operator[]} and \code{operator()}
of the basic classes call \code{laerror} (section~\ref{sec:concepts:errors}) when the reference count is larger
than one, since this is a possible write to shared data:
\begin{lstlisting}
NRVec<double> c = a;
c[1] = 7.;   //DEBUG: LA:ERROR - possible use of NRVec[] with count>1 as l-value
\end{lstlisting}
The check cannot distinguish reading from writing. Reading elements of a possibly shared object should
therefore be done through a const reference, which selects the const version of the operator:
\begin{lstlisting}
void f(NRVec<double> &x)   //non-const, x may share data with another vector
{
const NRVec<double> &xc = x;
double s = xc[0] + xc[1];  //no false alarm
x.copyonwrite();
x[0] = s;                  //write
}
\end{lstlisting}
The check can be switched off at run time by setting the global variable \code{LA::_LA_count_check=false}.
It is a safety net only: writes through raw pointers and writes in programs compiled without \code{DEBUG}
are not checked.

\lawarning{A missing \code{copyonwrite()} does not cause a crash, it silently modifies other objects which share
the data. It is the most frequent error when using the library. A library member function which modifies
data without calling \code{copyonwrite()} is a bug of the library.}

\subsection{Objects referring to external memory}

The constructor \code{NRVec(T *a, int n, bool skeleton)} with \code{skeleton=true} (and analogous
constructors of other classes, see the class chapters) creates an object which refers to memory owned by
somebody else, without copying. Its reference count is set to two, so that the memory is never deallocated by
the library and any modification by a library function first makes a private copy (writes through
\code{operator[]} go to the external memory, but are reported by the DEBUG check). It is useful for applying
library functions to data stored elsewhere. Such an object is intended as a \emph{read-only view} of the external data.

\section{Shallow and deep copy}\label{sec:concepts:assign}

\begin{lareference}
\lafunc{X(const X &rhs); X &operator=(const X &rhs)}{shallow copy: share the data with \code{rhs}
(section~\ref{sec:concepts:refcount}).}
\lafunc{X &operator|=(const X &rhs)}{deep copy: \code{*this = rhs} followed by \code{copyonwrite()}; this object gets its own copy
of the data, recursively for nested types.}
\lafunc{X &operator=(const T &a)}{assignment of a scalar. For vectors every element is set to \code{a}; for square matrices
(\code{NRMat}, \code{NRSMat}) the matrix is set to $a\,\mathbf{1}$, i.e.\ the diagonal is set to \code{a} and the
off-diagonal elements to zero. In the same spirit \code{+=} and \code{-=} with a scalar add to the diagonal of a matrix.}
\end{lareference}
The deep copy is needed when the copy will be modified by element access, or when a function
destroys its argument (e.g.\ \code{diagonalize}):
\begin{lstlisting}
NRMat<double> a(n,n);
...
NRMat<double> asave;
asave |= a;          //a will be overwritten
\end{lstlisting}
\code{x.swap(y)} exchanges the contents of two objects without any copying.

Conversions between different classes (e.g.\ \code{NRSMat} $\to$ \code{NRMat}, \code{NRMat} $\to$ \code{NRVec},
real $\to$ complex) are done by explicit constructors and always create new data, unless stated otherwise
in the class chapters.

\section{Traits and element types}\label{sec:concepts:traits}

\laheader{la_traits.h}

The algorithms of the library are written generically with the help of the traits class template
\code{LA_traits<T>}, which is defined for all scalar types and for all container classes of the library.
Its main members are:
\begin{lareference}
\lafunc{typedef ... elementtype}{type of the elements (for a scalar the type itself)}
\lafunc{typedef ... normtype}{type of the norm, e.g.\ \code{double} for \code{std::complex<double>} and for \code{NRMat<std::complex<double> >}}
\lafunc{typedef ... realtype; typedef ... complextype}{the corresponding real and complex type, e.g.\ \code{NRVec<double>} and \code{NRVec<std::complex<double> >}}
\lafunc{typedef ... producttype}{type of the product of two objects (for \code{NRSMat} it is \code{NRMat})}
\lafunc{static bool is_plaindata()}{true for scalar types, which can be copied by \code{memcpy}, false for the classes}
\lafunc{static bool is_complex()}{whether the (innermost) element type is complex}
\lafunc{static normtype norm(const T &x); static void axpy(T &s, const T &x, const elementtype c)}{norm and $s \leftarrow s + c\,x$, used by generic algorithms such as \code{matexp.h}}
\lafunc{static void copyonwrite(T &x); static void clearme(T &x); static void deallocate(T &x)}{generic versions of the member functions (no operation for scalars)}
\lafunc{static void put(int fd, const T &x, ...); static void get(int fd, T &x, ...); multiput; multiget}{binary I/O, recursive for nested types}
\lafunc{static T conjugate(const T &x); realpart(x); imagpart(x); setrealpart(x,y); setimagpart(x,y)}{complex helpers, trivial for real types}
\lafunc{static bool bigger(const T &x, const T &y); smaller(x,y)}{comparison used for sorting; for complex types it is an error}
\end{lareference}
The auxiliary template \code{isscalar<T>} distinguishes scalars from classes, \code{LA_traits_complex<T>}
gives for a complex type the component type and the corresponding real vector and matrix types (used by
the mixed real--complex functions such as \code{gemv} of a complex vector with a real matrix), and
\code{LA_traits_io<T>} makes \code{char} types be read and written as numbers by the stream operators.
The functions \code{LA::real(c)} and \code{LA::imag(c)} return references to the real and imaginary part of a
\code{std::complex}, which can be used as l-values.

\subsection{Element types and instantiation}

The class templates are explicitly instantiated in the library for \code{double}, \code{std::complex<double>}
and the integer types \code{char}, \code{short}, \code{int}, \code{long}, \code{long long} and their unsigned
variants. The BLAS and LAPACK backed functions are implemented for \code{double} and \code{std::complex<double>}
(some also for mixed real--complex arguments); calling them for other types results in a run-time
\code{laerror} (``not implemented for this type'') or a link error. Functions implemented in the headers work
for any type with the necessary arithmetic operations.

\subsection{Nested types}

The element type can be a class of the library, e.g.\ \code{NRVec<NRMat<double> >} (a list of matrices) or
\code{NRMat<NRVec<double> >}. Copy on demand is then applied on both levels: \code{copyonwrite()} of the outer
object copies the inner objects shallowly and calls their \code{copyonwrite()}, so that \code{|=} makes a
complete deep copy; \code{put} and \code{get} write and read the inner objects including their dimensions.
Arithmetic operations of the outer class (e.g.\ \code{axpy}) work elementwise via the traits.
\lawarning{Writing into an inner object obtained through \code{operator[]} of the outer object
(\code{v[i](j,k)=x}) requires both objects to be unshared: call \code{v.copyonwrite()} (which recursively
detaches the inner objects as well) before.}

\section{Storage conventions}\label{sec:concepts:storage}

\begin{itemize}
\item Indices run from 0. The classes \code{NRVec_from1}, \code{NRMat_from1} and \code{NRSMat_from1} are
	derived from the basic classes and differ only in their element access operators, which take indices from 1.
	They convert to and from their base classes without copying.
\item \code{NRVec}: contiguous array of \code{size()} elements.
\item \code{NRMat}: row major (C) order, element $(i,j)$ of an $n\times m$ matrix is at position $i\,m+j$;
	\code{a[i]} is a pointer to the $i$-th row. Note that BLAS and LAPACK use column major order, so the
	library internally works with the transposed matrices; the user does not need to care, except when calling
	Fortran routines directly.
\item \code{NRSMat}: lower triangle of a symmetric or hermitian matrix, packed by rows, element $(i,j)$ with
	$i \ge j$ is at position $i(i+1)/2+j$ (which is also the upper triangle packed by columns, the format of the
	BLAS/LAPACK packed routines). The macros/functions \code{SMat_index} compute these positions.
\item \code{Tensor}: the leftmost index changes fastest (chapter~\ref{ch:tensor}).
\end{itemize}
All classes can return a raw pointer to their data (e.g.\ \code{operator T*()} of \code{NRVec} and \code{NRMat}),
for passing them to external routines; the copy on demand rule for writes through pointers applies.

\section{GPU support (experimental)}\label{sec:concepts:gpu}

When compiled with \code{CUDALA} (configure finds CUBLAS), the data of \code{NRVec}, \code{NRMat} and
\code{NRSMat} can reside either in the main memory or in the memory of a GPU. The location is described by
the enumeration
\begin{lstlisting}
typedef enum {undefined=-1, cpu=0, gpu1=1, gpu2=2, gpu3=3, gpu4=4} GPUID;
\end{lstlisting}
Constructors take an optional location argument (default \code{undefined}, which means the global default
location \code{DEFAULT_LOC}, settable by \code{set_default_loc()}); \code{getlocation()} returns the location
and \code{moveto(loc)} moves the data. BLAS-like operations of objects on the GPU are performed by CUBLAS;
operands of a binary operation must reside in the same location. Operations which are not implemented for the
GPU call \code{laerror} (``Operation not implemented on GPU (yet). Use .moveto(0) first.''); element access
by \code{operator[]} and \code{operator()} is possible only in the main memory. Without \code{CUDALA},
\code{getlocation()} always returns \code{cpu} and \code{moveto()} does nothing, so that code written for
the GPU compiles also without it.

\lawarning{The GPU support is experimental and has not been thoroughly tested. It is not documented in
detail in this manual; the class chapters at most mention operations restricted to the CPU.}
The GPU code is currently not maintained: it uses the legacy CUBLAS API (\code{cublasInit}) and a compute capability
(\code{-arch sm_20}) not supported by current CUDA versions, and the \code{CUDALA} build has known compile errors. It will be revisited
with current NVIDIA hardware.

\section{Error handling and DEBUG checks}\label{sec:concepts:errors}

\laheader{laerror.h, laerror.cc}

Errors are reported by the macro \code{laerror(message)}, which calls \code{LA::laerror2(message, function)}
with the name of the calling function. \code{laerror2} prints the message and the function name to both
\code{std::cout} and \code{std::cerr}, the system error message if \code{errno} is set, the command line, the
user name, host name and working directory (and the list of files in the working directory if
\code{LA_error_list_files} is true), and then throws an exception of the class
\begin{lstlisting}
class LAerror : public std::exception {public: std::string msg; ... const char *what() const noexcept;};
\end{lstlisting}
which can be caught and printed by \code{operator<<}:
\begin{lstlisting}
try { ... }
catch(LAerror &e) {std::cerr << "error: " << e << std::endl;}
\end{lstlisting}
Errors detected by BLAS, CBLAS and ATLAS (via \code{xerbla} and its variants) are converted to \code{laerror}
as well. Users may call \code{laerror} in their own code. The function \code{sigtraceback(signal,flags)}
installs a handler printing a stack trace (by running \code{gdb} on the process) when the given signal arrives,
e.g.\ \code{sigtraceback(SIGSEGV,1)}; see chapter~\ref{ch:utilities}.

With \code{DEBUG} defined, the library checks index ranges, dimensions of operands, allocation of objects
and the reference count in element access (section~\ref{sec:concepts:refcount}). Without \code{DEBUG}, these
errors lead to undefined behaviour. Programs should be developed and tested with \code{DEBUG}.

\section{Input and output}\label{sec:concepts:io}

All classes support three kinds of I/O:
\begin{description}
\item[Stream operators] \code{operator<<} and \code{operator>>} write and read a human readable text format:
	the dimensions followed by the elements (a matrix by rows, one row per line). The input operator resizes the
	object. \code{char} elements are written as numbers.
\item[\code{fprintf}, \code{fscanf}] formatted output with a \code{printf} format for the elements and input
	from a C \code{FILE*} (see the class chapters).
\item[\code{put}, \code{get}] unformatted binary I/O to a POSIX file descriptor:
\begin{lstlisting}
void put(int fd, bool dimensions=1, bool transp=0, bool orca_format=false) const;
void get(int fd, bool dimensions=1, bool transp=0, bool orca_format=false);
\end{lstlisting}
	With \code{dimensions} true, a header with the dimensions is written before the data and \code{get} resizes the
	object accordingly; otherwise nothing but the elements is written and the object must already have the right size.
	The format is the binary representation of the machine (no conversion of endianity). Per class:
	\begin{itemize}
	\item \code{NRVec}: header $n$ and a padding \code{int} 1 (8-byte alignment), then the $n$ elements; \code{transp} is ignored.
	\item \code{NRMat}: header \code{nrows}, \code{ncols}, then the elements by rows. With \code{transp} the transposed matrix is
	written (dimensions swapped, elements column by column); \code{get} with \code{transp} reads such data back into the
	original orientation, i.e.\ it reads a file stored by columns.
	\item \code{NRSMat}: header $n$, $n$, then the packed lower triangle; \code{transp} is ignored.
	\item \code{SparseMat}: header \code{nrows}, \code{ncols} (\code{SPMatindex}) and the \code{int} flags \code{symmetric} and
	\code{nonzero}, then triples (row, column, element) and a terminating pair of indices $-1$; \code{SparseSMat}: header
	\code{nrows}, \code{ncols}, triples and the terminating pair. With \code{transp} rows and columns are swapped.
	\item \code{Tensor} and \code{fourindex} have their own \code{put}/\code{get} (see their chapters).
	\end{itemize}
	\code{orca_format} (dense classes only; \code{laerror} for the sparse ones) writes the header of the ORCA program: the
	dimensions followed by \code{sizeof(T)} instead of the padding; ORCA matrices are stored by columns, so they are read by
	\code{get(fd,true,true,true)}. Elements of class type (e.g.\ a vector of matrices) are written recursively by their own
	\code{put}, with the same \code{dimensions} flag; \code{transp} is passed on to them only by the transposed
	\code{NRMat} and the sparse matrices.
\end{description}
For storing many objects of the same size in a scratch file, see \code{AuxStorage} (chapter~\ref{ch:utilities}).
